\documentclass[aspectratio=169]{beamer}

\usetheme{Madrid}
\usecolortheme{default}
\usefonttheme{professionalfonts}

\usepackage{hyperref}

\hypersetup{
    colorlinks=true,
    linkcolor=blue,
    urlcolor=blue
}

\setbeamertemplate{navigation symbols}{}

\title[Combustion]{UCK 427E -- Combustion}
\subtitle{Introduction and Course Overview}
\author{Prof.\ Dr.\ Onur Tun\c{c}er}
\institute{Istanbul Technical University}
\date{February 12\textsuperscript{th}, 2026}

\begin{document}

\begin{frame}
  \titlepage
\end{frame}

\begin{frame}{Why Study Combustion?}
\begin{itemize}
  \item More than 80\% of global energy conversion involves combustion
  \item Central to aerospace propulsion systems:
  \begin{itemize}
    \item Gas turbines
    \item Rocket engines
    \item Afterburners
  \end{itemize}
  \item Critical for:
  \begin{itemize}
    \item Efficiency
    \item Emissions
    \item Stability and safety
  \end{itemize}
\end{itemize}
\end{frame}

\begin{frame}{Combustion Is Multidisciplinary}

\begin{block}{Core Disciplines}
\begin{itemize}
  \item Thermodynamics
  \item Fluid mechanics
  \item Chemical kinetics
  \item Heat and mass transfer
\end{itemize}
\end{block}

\pause

\begin{block}{Why It Is Challenging?}
\begin{itemize}
  \item Strongly coupled physics
  \item Nonlinear behavior
  \item Wide range of time and length scales
\end{itemize}
\end{block}
\end{frame}

\begin{frame}{How Combustion Is Studied Today}
\begin{itemize}
  \item Analytical models (limited but insightful)
  \item Experiments (expensive, complex)
  \item \alert{Computational modeling} (dominant tool today)
\end{itemize}

\vspace{0.5cm}

\begin{block}{This Course}

We combine:

\begin{itemize}
  \item Fundamental theory
  \item Computational modeling
  \item Physical interpretation
\end{itemize}
\end{block}
\end{frame}

\begin{frame}{Modeling Hierarchy in Combustion}

\begin{center}
\begin{tabular}{c}
CFD / LES / DNS \\ \hline
1-D Flames \\ \hline
0-D Reactors \\ \hline
Chemical Kinetics \& Thermodynamics
\end{tabular}
\end{center}

\vspace{0.4cm}

\begin{itemize}
  \item This course focuses on 0-D and 1-D models
  \item These form the foundation of all CFD combustion models
\end{itemize}
\end{frame}

\begin{frame}{Course Objectives}

\begin{itemize}
  \item Understand physical principles of combustion
  \item Analyze reacting flow systems
  \item Model combustion processes computationally
  \item Develop intuition for flame behavior and emissions
\end{itemize}
\end{frame}

\begin{frame}{What You Will Learn}

\begin{itemize}
  \item Chemical thermodynamics of reacting mixtures
  \item Chemical kinetics and reaction mechanisms
  \item Premixed and diffusion flames
  \item Spray and turbulent combustion (conceptual)
  \item Computational combustion modeling using C++
\end{itemize}

\end{frame}

\begin{frame}{Computational Tools}

\begin{block}{Software}
\begin{itemize}
  \item C++ (primary language)
  \item CANTERA (thermochemistry and kinetics)
  \item Python or MATLAB (optional post-processing)
\end{itemize}
\end{block}

\pause

\begin{block}{Why C++?}
\begin{itemize}
  \item Performance
  \item Industry relevance
  \item Full control over models
\end{itemize}
\end{block}
\end{frame}

\begin{frame}{Course Structure}
\begin{itemize}
  \item Lectures: theory and physical understanding
  \item Homeworks: computational modeling
  \item Midterm and final exam
\end{itemize}

\vspace{0.3cm}

\begin{block}{Expectation}
This is a senior-level engineering course requiring:
\begin{itemize}
  \item Mathematical maturity
  \item Programming discipline
\end{itemize}
\end{block}
\end{frame}

\begin{frame}{Homework Philosophy}
\begin{itemize}
  \item Homework assignments are computational
  \item Focus on understanding, not coding volume
  \item Clean, readable, reproducible code expected
\end{itemize}

\vspace{0.3cm}

\begin{block}{Tentative Topics}
\begin{itemize}
  \item Adiabatic flame temperature
  \item Auto-ignition
  \item Laminar flame speed
  \item Emissions formation
\end{itemize}
\end{block}
\end{frame}

\begin{frame}{Evaluation}
\begin{center}
\begin{tabular}{l c}
Midterm Exam & 10\% \\
Homeworks (5) & 50\% \\
Final Exam & 40\%
\end{tabular}
\end{center}
\end{frame}

\begin{frame}{What Comes Next?}
\begin{itemize}
  \item Review of thermochemistry
  \item Ideal gas mixtures
  \item Enthalpy, entropy, and equilibrium
\end{itemize}

\vspace{0.5cm}

\centering
\alert{Please install Cantera before next week}
\end{frame}

\begin{frame}{Cantera Resources}

\begin{block}{Official Website}
\centering
\Large
\href{https://cantera.org}{\textbf{cantera.org}}
\end{block}

\vspace{0.4cm}

\begin{itemize}
  \item Documentation
  \item C++ Tutorials
  \item Examples and YAML mechanisms
\end{itemize}

\end{frame}



\begin{frame}{Thermodynamic Description of Matter}
\begin{block}{Fundamental Question}
How do we describe the state of a reacting system?
\end{block}

\vspace{0.3cm}

A thermodynamic state is defined by a set of measurable properties.

\vspace{0.3cm}

Examples:
\begin{itemize}
  \item Pressure
  \item Temperature
  \item Volume
  \item Internal energy
  \item Enthalpy
  \item Entropy
\end{itemize}
\end{frame}

\begin{frame}{Intrinsic vs Extrinsic Properties}

\begin{block}{Intrinsic (Intensive) Properties}
Independent of system size or mass.
\end{block}

Examples:
\begin{itemize}
  \item Temperature
  \item Pressure
  \item Density
  \item Specific internal energy
  \item Specific enthalpy
\end{itemize}

\pause

\begin{block}{Extrinsic (Extensive) Properties}
Depend on the size or amount of matter.
\end{block}

Examples:
\begin{itemize}
  \item Mass
  \item Total internal energy
  \item Total enthalpy
  \item Volume
  \item Entropy
\end{itemize}

\end{frame}

\begin{frame}{Mathematical Definition}

\begin{block}{Extensive Property}
If system size is scaled by factor $\lambda$:

\[
X \rightarrow \lambda X
\]

\end{block}


\begin{block}{Intensive Property}
Remains unchanged when system size is scaled:

\[
x = \frac{X}{m}
\]

\end{block}

\vspace{0.3cm}

Example:

\[
u = \frac{U}{m}, \quad h = \frac{H}{m}
\]

\end{frame}

\begin{frame}{Why This Matters in Combustion}

\begin{itemize}
  \item Combustion chambers contain reacting mixtures
  \item Chemical reactions change composition
  \item We must distinguish:
  \begin{itemize}
    \item Total energy in chamber (extensive)
    \item Energy per unit mass (intensive)
  \end{itemize}
\end{itemize}

\vspace{0.3cm}

\begin{block}{Key Idea}
Reaction rates depend on intensive properties.
\end{block}

\end{frame}

\begin{frame}{Mixtures and Specific Quantities}

For a reacting mixture:

\[
U = \sum_i m_i u_i
\]

Specific internal energy:

\[
u = \sum_i Y_i u_i
\]

where:
\begin{itemize}
  \item $Y_i$ = mass fraction
  \item $u_i$ = specific internal energy of species $i$
\end{itemize}

\end{frame}

\begin{frame}{State Postulate for Reacting Systems}

For a simple compressible reacting system:

\[
\text{State} = f(T, P, \{Y_i\})
\]

\vspace{0.5cm}

Temperature + Pressure + Composition
\end{frame}

\begin{frame}
\centering
\Large Questions?
\end{frame}

\end{document}

